\documentclass[11pt]{article}
\usepackage{fontspec,amsmath,amssymb,amsthm,hyperref}
\usepackage[margin=1in]{geometry}
\setlength{\emergencystretch}{3em}
\title{Construct the mean-field L\ensuremath{{}^{2}} integral equations}
\author{LangevinLean mathematical exposition}
\date{}
\begin{document}\maketitle
\noindent Source module: \texttt{ProximalLangevin/MeanFieldProcess.lean}. Compile this proof document with XeLaTeX or LuaLaTeX. The accompanying HTML pairs the mathematical argument with the exact Lean statements and records the scope of the formalization.


Work on a probability space \(Q\), with \(H=L^2(Q;\mathbb R^{d_\theta+d_x})\). A potential \(V\) with globally \(L\)-Lipschitz gradient defines the actual mean-field drift. Its continuous finite-horizon integral solution is constructed by the Banach-space Picard theorem.


\section{Expectation is a bounded linear map from L\ensuremath{{}^{2}}}


The norm-of-integral inequality and \(\|X\|_{L^1}\le\|X\|_{L^2}\) on a probability space give \(\|\mathbb E_Q X\|\le\|X\|_2\). Bochner additivity and scalar linearity therefore define a bounded linear map \(\mathcal E:H\to E\) with operator norm at most one. The evaluation lemma records that its value is the actual integral.




\textbf{Correspondence with the displayed declarations.}


\begin{itemize}

\item \texttt{norm\_lp\_expectation\_le}: On a probability space, the actual expectation is bounded by the L\ensuremath{{}^{2}} norm.
\item \texttt{lpExpectationCLM}: Actual Bochner expectation as a bounded linear map out of L\ensuremath{{}^{2}}.
\end{itemize}


\noindent Lean declarations: \texttt{norm\_lp\_expectation\_le}, \texttt{lpExpectationCLM}, \texttt{lpExpectationCLM\_apply}, \texttt{lpExpectationCLM\_norm\_le}.

\section{Coordinate projections and a parameter inclusion}


Define \(P_x(\theta,x)=x\) and \(I_\theta(\theta)=(\theta,0)\). Their coordinate identities follow by evaluating the definitions. The Euclidean square identity \(\|(\theta,0)\|^2=\|\theta\|^2\), followed by nonnegative square-root cancellation, proves that \(I_\theta\) is an isometry.




\textbf{Correspondence with the displayed declarations.}


\begin{itemize}

\item \texttt{jointLatentCLM}: The latent component of one parameter/latent joint vector.
\item \texttt{jointParameterInclusionCLM}: Embed only a parameter block into the joint space.
\end{itemize}


\noindent Lean declarations: \texttt{jointLatentCLM}, \texttt{jointLatentCLM\_apply}, \texttt{jointParameterInclusionCLM}, \texttt{jointParameterInclusionCLM\_parameter}, \texttt{jointParameterInclusionCLM\_latent}, \texttt{jointParameterInclusionCLM\_norm}.

\section{A true latent mask erases all parameter coordinates}


Set \(M_x=I-I_\theta P_\theta\). Its parameter block is zero and its latent block is unchanged. Expanding the Euclidean squared norm shows \(\|M_xz\|^2\le\|z\|^2\), since the discarded parameter square is nonnegative. Thus the mask is nonexpansive and remains bounded after composition on \(L^2\).




\textbf{Correspondence with the displayed declarations.}


\begin{itemize}

\item \texttt{jointLatentMaskCLM}: The actual latent-coordinate projection on the joint space.
\end{itemize}


\noindent Lean declarations: \texttt{jointLatentMaskCLM}, \texttt{jointLatentMaskCLM\_parameter}, \texttt{jointLatentMaskCLM\_latent}, \texttt{norm\_jointLatentMaskCLM\_le}.

\section{Average exactly the parameter block}


Construct the bounded linear map

\[\mathcal M X=\operatorname{const}(I_\theta P_\theta\mathbb E_QX)+M_xX.\]

The expectation estimate and coordinate contractions bound each term by \(\|X\|_2\), so \(\|\mathcal M X\|_2\le2\|X\|_2\). Linearity then proves \(\operatorname{Lip}(\mathcal M)\le2\). This is a conservative sufficient constant.




\textbf{Correspondence with the displayed declarations.}


\begin{itemize}

\item \texttt{meanFieldAverageCLM}: Average the parameter coordinates and retain the latent coordinates as actual random variables. This is a bounded linear map, not an assumed drift.
\end{itemize}


\noindent Lean declarations: \texttt{meanFieldAverageCLM}, \texttt{meanFieldAverageCLM\_norm\_le}, \texttt{meanFieldAverageCLM\_lipschitz}.

\section{Define the actual drift and apply Picard}


Set \(\mathcal A_V(X)=\mathcal M(\nabla V(X))\). The Nemytskii map for the actual Lipschitz gradient preserves \(L^2\) and has constant \(L\), so \(\operatorname{Lip}(\mathcal A_V)\le2L\). For any continuous \(H\)-valued input \(B\) with \(B_0=0\), and any \(T\ge0\), Picard constructs continuous \(Y\) with \(Y_0=Y^{\rm in}\) and

\[Y_t-Y^{\rm in}=-\int_0^t\mathcal A_V(Y_s)ds+B_t,\qquad0\le t\le T.\]

The proof is an existence theorem in the actual complete Banach space \(H\).




The averaging operator retains the latent gradient pointwise and replaces the parameter gradient by its expectation as a constant random variable. Expectation is a bounded linear map \(L^2\to H\) of norm at most one on a probability space, and the coordinate masks are orthogonal contractions. Thus the averaged drift is globally Lipschitz once the true joint score is Lipschitz. Picard's theorem applies in the complete Hilbert space of square-integrable joint states to a continuous additive input, constructing the actual integral solution on every finite horizon.




\textbf{Correspondence with the displayed declarations.}


\begin{itemize}

\item \texttt{lpMeanFieldDrift}: The actual mean-field drift: expected parameter force and pointwise latent force.
\item \texttt{lpMeanFieldDrift\_lipschitz}: A conservative global Lipschitz constant for the constructed drift.
\item \texttt{exists\_lp\_meanField\_integral}: Picard constructs the continuous mean-field L\ensuremath{{}^{2}} integral solution from its actual gradient and a continuous L\ensuremath{{}^{2}} additive input on every finite horizon.
\end{itemize}


\noindent Lean declarations: \texttt{lpMeanFieldDrift}, \texttt{lpMeanFieldDrift\_lipschitz}, \texttt{exists\_lp\_meanField\_integral}.

\section{Verify the intended block formulas in L\ensuremath{{}^{2}}}


Almost everywhere representative equalities show

\[(\mathcal M X)(\omega)=(\mathbb E_Q P_\theta X,\,P_xX(\omega)).\]

Applying parameter projection gives the constant random variable \(\operatorname{const}(P_\theta\mathbb E_QX)\); applying latent projection gives \(P_xX\). The proofs explicitly combine the almost everywhere addition, constant and composition identities, so they are equalities in \(L^2\).




\textbf{Correspondence with the displayed declarations.}


\begin{itemize}

\item \texttt{meanFieldAverageCLM\_coe}: The averaging map has exactly its intended pointwise block action.
\item \texttt{meanFieldAverageCLM\_parameter}: Averaging makes the parameter block a constant random variable.
\item \texttt{meanFieldAverageCLM\_latent}: Averaging leaves the actual latent random variable unchanged.
\end{itemize}


\noindent Lean declarations: \texttt{meanFieldAverageCLM\_coe}, \texttt{meanFieldAverageCLM\_parameter}, \texttt{meanFieldAverageCLM\_latent}.

\section{Identify the parameter force with the true marginal integral}


Define \(a_V(X)=P_\theta\mathbb E_Q\nabla V(X)\), a continuous map on \(H\). The drift blocks are \(\operatorname{const}(a_V(X))\) and \(P_x\nabla V(X)\). If \(P_\theta X=\operatorname{const}(\theta)\), reconstruct \(X=(\theta,P_xX)\) almost everywhere. Commute the bounded projection through expectation and use the pushforward integral formula to prove

\[a_V(X)=\int\nabla_\theta V(\theta,x)\,\mathcal L_Q(P_xX)(dx).\]


\textbf{Correspondence with the displayed declarations.}


\begin{itemize}

\item \texttt{lpMeanFieldParameterForce}: The actual expected parameter force of a square-integrable joint vector.
\item \texttt{lpMeanFieldDrift\_parameter}: The parameter projection of the actual drift is its expected force.
\item \texttt{lpMeanFieldDrift\_latent}: The latent projection of the actual drift is pointwise latent gradient.
\item \texttt{lpMeanFieldParameterForce\_eq\_meanFieldDrift}: The expected force equals the actual mean-field drift integral whenever the input parameter is deterministic. The latent marginal is its actual law.
\end{itemize}


\noindent Lean declarations: \texttt{lpMeanFieldParameterForce}, \texttt{lpMeanFieldParameterForce\_continuous}, \texttt{lpMeanFieldDrift\_parameter}, \texttt{lpMeanFieldDrift\_latent}, \texttt{lpMeanFieldParameterForce\_eq\_meanFieldDrift}.

\section{Project the integral equation and construct a deterministic parameter}


Suppose the initial parameter is constant \(\theta_0\) and \(P_\theta B_t=0\). Bounded linear maps commute with the Bochner time integral, so projecting the constructed equation gives \(P_\theta Y_t=\operatorname{const}(\theta_t)\), where \(\theta_t=\theta_0-\int_0^t a_V(Y_s)ds\). The continuous primitive proves continuity of \(\theta\). Set \(\mu_t=\mathcal L_Q(P_xY_t)\) and insert the force identity:

\[\theta_t-\theta_0=-\int_0^t b_V(\theta_s,\mu_s)ds,\qquad P_xY_t-P_xY_0=-\int_0^tP_x\nabla V(Y_s)ds+P_xB_t.\]

Both identities hold in their stated deterministic or \(L^2\) sense on the finite horizon.




Project the constructed joint equation onto the parameter and latent blocks. Continuous linear maps commute with the Bochner time integral. Since the additive input has zero parameter block, its parameter coordinate remains a constant random variable, represented by a deterministic continuous path \(\theta_t\). The force identity replaces expectation of the joint parameter score by its integral under the actual latent pushforward law \(\mu_t\), giving \[\theta_t-\theta_0=-\int_0^t\left[\int\nabla_\theta V(\theta_s,x)\,d\mu_s(x)\right]ds.\] The latent projection yields the corresponding drift-plus-noise equation. The parameter law is therefore derived from the constructed process, not introduced as a free auxiliary measure path.




\textbf{Correspondence with the displayed declarations.}


\begin{itemize}

\item \texttt{lp\_meanField\_parameter\_equation}: Projection of the constructed equation preserves a deterministic parameter and gives its genuine integral equation.
\item \texttt{lp\_meanField\_latent\_equation}: The latent component satisfies the actual projected gradient equation.
\item \texttt{exists\_lp\_meanField\_equations}: The constructed solution has a continuous deterministic parameter, its actual latent marginal, and both mean-field integral equations. No solution trajectory or law evolution is assumed in this existence statement.
\end{itemize}


\noindent Lean declarations: \texttt{lp\_meanField\_parameter\_equation}, \texttt{lp\_meanField\_latent\_equation}, \texttt{exists\_lp\_meanField\_equations}.

\section{Construct a continuous latent-only Brownian input}


Apply the latent mask to the actual coordinate Brownian \(L^2\) noise. Brownian second moments prove its continuity in \(L^2\) and its zero value at time zero. Masking makes its parameter projection zero at every time. No separate continuity assumption on a stochastic error process is used.




\textbf{Correspondence with the displayed declarations.}


\begin{itemize}

\item \texttt{meanFieldBrownianNoise}: Actual Brownian-coordinate inputs, with only their latent coordinates retained, give a concrete continuous L\ensuremath{{}^{2}} input. Noise continuity and its zero initial value are conclusions of the Brownian laws.
\end{itemize}


\noindent Lean declarations: \texttt{meanFieldBrownianNoise}, \texttt{meanFieldBrownianNoise\_continuous}, \texttt{meanFieldBrownianNoise\_zero}, \texttt{meanFieldBrownianNoise\_parameter}.

\section{Specialize to Brownian input and to the split envelope}


The constructed Brownian input satisfies all hypotheses of the general integral existence theorem, so substitution gives the Brownian-driven \(L^2\) solution. For \(V=g_1+g_2^\eta\), the actual proximal calculus supplies the Lipschitz gradient constant \(L_1+\eta^{-1}\), which is included as a derived witness in the conclusion. These results establish continuous \(L^2\) integral equations; they do not assert a pathwise strong-SDE or a Fokker–Planck PDE theorem.




\textbf{Correspondence with the displayed declarations.}


\begin{itemize}

\item \texttt{exists\_lp\_meanField\_integral\_of\_brownian}: Genuine finite-horizon mean-field L\ensuremath{{}^{2}} existence with a Brownian input. All stochastic input continuity is derived from the coordinate Brownian laws; no pathwise strong-SDE, independence or adaptedness assertion is made.
\item \texttt{exists\_lp\_moreau\_meanField\_integral}: Split Moreau smoothing supplies the actual globally Lipschitz gradient needed by the mean-field construction. Its coefficient is \texttt{L\ensuremath{{}_{1}}+1/eta}.
\end{itemize}


\noindent Lean declarations: \texttt{exists\_lp\_meanField\_integral\_of\_brownian}, \texttt{exists\_lp\_moreau\_meanField\_integral}.

\end{document}
